\section{Limitations}\label{limitation}

The findings above are qualified by several of limitations.
Some concerns what estimates can be taken to mean and others follow from the size of the sample.
Exchange rate volatility and economic growth are jointly determined.
A currency may become volatile because the economy is weakening or just as a volatile currency may weaken the economy.
This design cannot tell which is happening.
A negative coefficient is consistent.
Janus and Riera-Crichton (\citeyear{thorsten_pdf_2015}) address with this finding a variable that moves exchange rate volatility but has no direct effect on growth.
Their case using volatility of a country's commodity terms of trade and isolate the part of volatility that cannot be caused by growth.
There are no similar variables in this analysis.
The estimates in this thesis are association instead of the causal effects.
The volatility is estimated and not observed.
Volatility is not measured directly but estimated in a separate step.
The uncertainty in that estimate is not passed on to the panel regression.
The regression treats the annual volatility average as they were observed data when each carries its own estimation error.
The result is significance overstated by an unknown margin.
Bootstrapping the two steps together would address this.
But it is not performed here.

\vspace{\baselineskip}

The sample is small in both dimensions.
Six countries and twenty years is not much data to work with and it affects the standard errors.
Driscoll-Kraay standard errors reply on having many time periods instead of many countries.
Twenty years is at the low end of what the method needs.
The estimator is consistent and it converges as the sample grows but whether well at this length of panel is a different question.
So the p values presented should be read as approximate instead of exact values.

\vspace{\baselineskip}

The sample cannot support strong claims about stability.
With six countries removing one country means removed a sixth of the observations.
And result will be sensitive to which countries are included.
The crisis years are similar because they are the only period which volatility differs heavily across the sample.
Both problems cannot be fix within this design.
The six countries at two ends of the range with nothing between two groups.
Any estimate of how the effect changes across the range is on the linear form assumed by the specification instead of an observed variation.
The model cannot detect the problem in the relationship where it has no observations.
This follow from which currencies kept their own currencies and was not the study could arranged differently.

\vspace{\baselineskip}

There is limitation on how the variables are measured.
The volatility is measured in nominal terms.
Exchange rates against the euro are used as published without adjusting for relative price levels.
This suits the mechanism the thesis examines since a firm contracting in euro faces the nominal rate but it means the measure cannot separate currency movement from price stability.
A country with unstable inflation will show high nominal volatility.
And the coefficient cannot tell the two apart.
The correlation between volatility and inflation in this sample at 0.29 so the problem present.
A real effective rate would avoid it but it is measuring the movement against the world instead of the euro.
Also GARCH cannot tell a pegged currency.
The model works by describing how volatility changes over time which requires volatility that can changes.
A central bank holds currency close to fixed rate so there is nothing for the model to describe.
This applies to Denmark in the whole sample period in weaker form.
Also Romania under its managed float and Czech Republic exchange rate floor regime during 2013 to 2017.
The estimates still capture how much each currency moved but for three of the six countries they do not describe a volatility process.
